The main evidence supports a focused claim: parameterized simulation stress tests can reveal target-generation failures that are hidden by clean-setting success rates. The result is strongest for diagnostic target generation and detector transfer, not for end-to-end closed-loop manipulation. This distinction matters because the benchmark is designed to isolate the language-to-target stage before policy execution.

The threshold analysis changes the interpretation of the RGB-D baselines. Crop-median depth should not be described as universally strongest. It is stronger at the default and relaxed thresholds, while box-center depth is more precise under strict thresholds. This precision--robustness tradeoff is useful because it tells benchmark users which target sources are likely to fail when tolerances change.

The open-vocabulary bridge also narrows the language-conditioned claim. Metadata-assisted query-aware selection is a controlled comparator, first detection is a stress-sensitive comparator, and GroundingDINO is a real detector plug-in tested through the same interface. The YCB/clutter held-out result is negative but informative: the generic detector query and current adapter labels do not transfer to the harder YCB/clutter setting in this setup. The query ablation further shows that changing generic prompts to object-label or category templates does not resolve the failure because the current adapter exposes generic YCB labels. A reproducible benchmark should expose such detector-query and adapter-metadata failures explicitly rather than masking them behind metadata-only evaluation.

The closed-loop smoke result is the largest remaining limitation. Although diagnostic success is useful for target-generation analysis, the present scripted executor did not pass the oracle task-success gate. This prevents a positive claim that target-distance success is already calibrated to actual task success. Future work should add a validated scripted or learned executor, then report agreement, target-error bins, and task-success prediction under the same seed splits.

Other limitations are also deliberate. The study is simulation-only; it does not claim real-robot robustness. It does not evaluate state-of-the-art VLA policies such as RT-style or OpenVLA-style controllers. It does not claim that GroundingDINO is the best open-vocabulary detector or that the current prompt configuration is optimal. These boundaries keep the contribution centered on reproducible diagnosis rather than on policy benchmarking or detector leaderboards.
